\clearpage
\section{Coefficients available in the manuscript}
\label{sec:coefficients}
We separate an explicit result in the supplied manuscript~\cite{theory}
from a coefficient estimated by fitting finite-chain data. Throughout this
section, $\alpha$ denotes the leading power-law exponent. Every reported
coefficient $c_r$ multiplies $x^r$ itself and includes all powers of $\pi$.
For example, the boson vacuum--spin coefficient is $\pi^2/12$, not
$1/12$. The same convention applies to $y^\beta$ for large intervals.
These expansions concern the continuum,
small-interval limit; they are not exact formulas at fixed two-site or
three-site intervals.

The previous report added derivations of boson $x^4,x^6$ coefficients and
Ising spin-channel $x^4$ coefficients. Those were additional calculations,
not coefficients stated in the original paper. They are excluded from the
current theoretical comparisons and preserved separately in the task's
archive. No missing subleading coefficient is derived in this revision.

\subsection{Ising: two explicit subleading coefficients}
At $p=y=1/2$, Eqs.~(4.15)--(4.16) of the supplied manuscript~\cite{theory} give
\begin{align}
 D_{I\mid\eps}
 &=\frac{2\pi^4}{15}x^4-\frac{8\pi^6}{45}x^6+\OO(x^8),
 \label{eq:Ieps}\\
 D_{\eps\mid I}
 &=\frac{2\pi^4}{15}x^4-\frac{8\pi^6}{21}x^6+\OO(x^8).
 \label{eq:epsI}
\end{align}
Here $D_{a\mid b}=S(\rho_A^a\|(\rho_A^a+\rho_A^b)/2)$, where
$\rho_A^a$ is the reduced density matrix of state $a$ on the interval $A$.
The substitution into the paper's general-weight expressions is
$(8/15)(1/4)=2/15$ for the leading coefficient and
$(32/315)(1/4)(8-15)=-8/45$ or
$(32/315)(1/4)(8-23)=-8/21$ for the subleading coefficient.
Thus the coefficients of $x^6$ itself are $-8\pi^6/45$ and
$-8\pi^6/21$, respectively.

For each of $I\mid\sigma$, $\sigma\mid I$, $\sigma\mid\eps$, and
$\eps\mid\sigma$, Eqs.~(4.13), (4.14), (4.17), and (4.18)
of the supplied manuscript~\cite{theory} give only
\begin{equation}
 D_{a\mid b}=\frac{\pi^2}{48}x^2+\OO(x^4).
 \label{eq:spinleading}
\end{equation}
The paper explicitly leaves the fourth-order four-energy-field
contribution unevaluated. Its discussion identifies additional sectors
at this order, but does not give a complete $x^4$ coefficient for these
four directions. The shared leading coefficient does not assert equality
of the complete directed relative entropies.

\begin{table}[htbp]\centering
\begin{tabular}{lrrrl}\toprule
Direction & $\alpha$ & $c_\alpha$ & $c_{\alpha+2}$ & Supplied equation\\\midrule
$I\mid\sigma$ &2&$\pi^2/48$&---&(4.13)\\
$\sigma\mid I$ &2&$\pi^2/48$&---&(4.14)\\
$I\mid\eps$ &4&$2\pi^4/15$&$-8\pi^6/45$&(4.15)\\
$\eps\mid I$ &4&$2\pi^4/15$&$-8\pi^6/21$&(4.16)\\
$\sigma\mid\eps$ &2&$\pi^2/48$&---&(4.17)\\
$\eps\mid\sigma$ &2&$\pi^2/48$&---&(4.18)\\\bottomrule
\end{tabular}
\caption{Complete coefficients supplied for the six Ising directions at
$p=1/2$, in powers of $x$. A dash means that the complete coefficient
is not supplied; it does not mean zero. Equation numbers in the last column
refer to the supplied manuscript~\cite{theory}. Higher coefficients beyond those
shown are not supplied for these directed mixture formulas.}
\label{tab:isingcoeff}
\end{table}

\subsection{Compact boson: leading coefficients only}
Write the left and right charge vector as $\boldsymbol Q_a$:
$\boldsymbol Q_{00}=(0,0)$,
$\boldsymbol Q_{ss}=(1/\sqrt2,1/\sqrt2)$, and
$\boldsymbol Q_{3s,s}=(3/\sqrt2,1/\sqrt2)$.
These are the charges denoted by $(Q_L,Q_R)$ in
Sec.~\ref{sec:bosonconstruction}; $\alpha$ denotes the leading exponent.
Define the squared charge separation
$\kappa_{ab}=\|\boldsymbol Q_a-\boldsymbol Q_b\|^2$.
Applying the paper's leading current-sector formula at $p=1/2$ gives
\begin{equation}
 \Dab=\frac{\kappa_{ab}\pi^2}{12}x^2+\OO(x^4).
 \label{eq:bosonresult}
\end{equation}
The original manuscript does not provide complete $x^4$ or $x^6$
coefficients for these six equal-mixture directions. General formulas
for individual subleading sectors do not supply the sum of all
contributions at the same order.

\begin{table}[htbp]\centering
\begin{tabular}{lrrr}\toprule
Directions, each separately & $\kappa_{ab}$ & $\alpha$ & $c_\alpha$\\\midrule
$00\mid ss,\quad ss\mid00$ &1&2&$\pi^2/12$\\
$00\mid3s,s,\quad3s,s\mid00$ &5&2&$5\pi^2/12$\\
$ss\mid3s,s,\quad3s,s\mid ss$ &2&2&$\pi^2/6$\\\bottomrule
\end{tabular}
\caption{Boson leading coefficients at $p=1/2$ in powers of $x$.
Each direction has an $\OO(x^4)$ remainder whose complete coefficient
is not supplied. Equal leading coefficients do not constrain
finite-chain reverse directions to be equal.}
\label{tab:bosoncoeff}
\end{table}

\subsection{Unknown coefficients in the fitting ansatz}
To test corrections numerically without deriving their coefficients,
we use only the leading and the next remainder power identified here:
\begin{equation}
 F_1(x)=c_0x^\alpha,\qquad F_2(x)=c_0x^\alpha+c_1x^{\alpha+2},
 \qquad
 \alpha=\begin{cases}
 4,&\{a,b\}=\{I,\eps\},\\
 2,&\text{all other requested pairs}.
 \end{cases}
 \label{eq:hierarchy}
\end{equation}
The subscripts 0 and 1 label the two coefficients, not their exponents.
The paper fixes the leading exponent and the next remainder order for
these cases. A remainder does not establish a nonzero coefficient, and
we do not assume an infinite sequence of even powers. Both coefficients
in $F_2$ are fitted. Section~\ref{sec:fits} additionally fits a single
power with its exponent free, to test whether the data recover $\alpha$.
